\section{Methodology}

In this section, we present the overall architecture of the proposed Unified Dual-Mask Physical Model, termed GeometricDualMask, for non-Lambertian light field depth estimation. Unlike previous works that predominantly rely on the idealized Lambertian reflectance assumption [7], our architecture explicitly addresses the physical boundary where Epipolar Plane Image (EPI) [7] linear structures degrade due to complex material reflections. As illustrated in Fig. 1, the framework integrates a physics-driven signal decomposition mechanism with a dual-branch geometric routing network. By decoupling the angular signals into illumination, disparity, and material components, the model adaptively processes Lambertian and non-Lambertian regions through specialized branches, ultimately yielding a unified and highly accurate depth representation.

The network takes light field data as input, comprising a $9 \times 9$ angular sampling grid and the corresponding center view image. Let the spatial dimensions of the center view be $H \times W$. The angular patch for each spatial pixel $(x, y)$ is denoted as $I(u,v) \in \mathbb{R}^{9 \times 9 \times C}$, where $(u,v)$ represents the angular coordinates and $C$ is the number of color channels. Prior to entering the main network, the raw light field data undergoes standard preprocessing, including normalization and the extraction of 2D EPI slices along the horizontal and vertical angular dimensions. These EPI slices, alongside the center view, serve as the foundational geometric priors for subsequent feature extraction, ensuring that the intrinsic epipolar geometry is preserved for the network to exploit.

The data flow progresses through four primary modules, each designed to tackle specific physical characteristics of light field imaging. First, the \textbf{Three-Layer Angular Signal Decomposition} module physically decouples the discrete angular sampling signals. Motivated by the need to separate material-induced reflections from geometric disparities, this module decomposes $I(u,v)$ into a direct current (DC) component representing ambient illumination, linear coefficients $a$ and $b$ capturing disparity-induced EPI slopes, and a residual component $\varepsilon(u,v)$ encoding material-specific Bidirectional Reflectance Distribution Function (BRDF) [7] properties. The decomposition is formulated as:
\begin{equation}
\label{eq:eq1}
I(u,v) = DC + a \cdot u + b \cdot v + \varepsilon(u,v)
\end{equation}
Following the physical decoupling, angular gradients are further extracted as geometric pseudo-labels to guide the subsequent mask generation. This explicit physical decoupling prevents the entanglement of texture and depth cues, a common limitation in end-to-end learning approaches. Following decomposition, the features are routed into a dual-branch architecture. The \textbf{Lambertian EPI Branch} processes the EPI features, linear coefficients, and the center view. Since Lambertian surfaces satisfy the angular cone constraint, this branch employs convolutional layers to extract linear EPI structures and incorporates a continuous depth prior, such as the Plane Regular Sampling Operator (PRSO) [7], to enforce piecewise smoothness. The depth prediction for Lambertian regions, $D_L \in \mathbb{R}^{H \times W \times C_{epi}}$, is regressed as:
\begin{equation}
\label{eq:eq2}
D_L = \text{PRSO}(\text{Conv}(\text{EPI\_Features}))
\end{equation}
Conversely, the \textbf{Non-Lambertian Geometric Branch} handles regions where EPI assumptions fail, such as specular highlights and scattering. Instead of relying on corrupted frequency or linear features, this branch utilizes the residual $\varepsilon(u,v)$ and angular gradients to capture local geometric anomalies (e.g., X-shaped patterns in EPIs). The non-Lambertian depth prediction, $D_{NL} \in \mathbb{R}^{H \times W \times C_{geo}}$, is derived via:
\begin{equation}
\label{eq:eq3}
D_{NL} = \text{Conv}(\text{Geometric\_Features}(\varepsilon, \text{angular\_gradient}))
\end{equation}
This specialized routing ensures that non-Lambertian regions are processed with appropriate geometric cues rather than being forced into an incompatible linear EPI model. Finally, the \textbf{Dual-Mask Generation and Fusion} module integrates the dual-branch outputs. To address the severe channel imbalance between the two branches (e.g., $C_{epi}=320$ vs. $C_{geo}=3$), both $D_L$ and $D_{NL}$ undergo a channel-balanced linear projection to a unified dimension ($\text{FUSE\_DIM}=96$), yielding projected features $F_L$ and $F_{NL}$. The selection of $\text{FUSE\_DIM}=96$ is empirically justified through ablation studies, as it achieves an optimal trade-off between representational capacity and computational overhead, preventing the high-dimensional EPI features from overshadowing the low-dimensional geometric cues during fusion. The projected features are concatenated and passed through convolutional layers to generate a pixel-level domain classification mask $M \in \mathbb{R}^{H \times W \times 1}$, which distinguishes Lambertian from non-Lambertian regions:
\begin{equation}
\label{eq:eq4}
M = \sigma(\text{Conv}(\text{Proj}(F_L) \oplus \text{Proj}(F_{NL})))
\end{equation}
where $\sigma$ denotes the sigmoid activation and $\oplus$ represents channel-wise concatenation. The final unified depth map $D_{final}$ is obtained by mask-weighted fusion:
\begin{equation}
\label{eq:eq5}
D_{final} = M \odot D_L + (1-M) \odot D_{NL}
\end{equation}
where $\odot$ indicates element-wise multiplication. This fusion strategy dynamically adapts to mixed-material scenes, significantly enhancing depth accuracy across diverse surface types.

The architecture outputs two primary tensors: the unified high-precision depth map $D_{final} \in \mathbb{R}^{H \times W \times 1}$ and the pixel-level domain classification mask $M \in \mathbb{R}^{H \times W \times 1}$. During training, the mask $M$ is supervised using the extracted angular gradient pseudo-labels via a domain classification loss ($\mathcal{L}_{domain}$), replacing less effective scene-level binary cross-entropy supervision. Simultaneously, the depth map is optimized using a depth regression loss ($\mathcal{L}_{depth}$). The total objective function is formulated as $\mathcal{L}_{total} = \lambda_d \mathcal{L}_{depth} + \lambda_m \mathcal{L}_{domain}$, where the loss weights $\lambda_d$ and $\lambda_m$ are both set to 1.0 to balance the optimization. In the post-processing and training sampling stage, a weighted random sampler and domain-balanced sampling strategy are employed to maintain multi-domain exposure balance. Specifically, the domain-balanced sampling strategy dynamically adjusts the sampling probabilities based on the proportion of non-Lambertian pixels in each batch. This mechanism specifically oversamples scarce non-Lambertian training instances (with an oversampling factor of 10), ensuring robust generalization and preventing the network from biasing towards dominant Lambertian regions in complex mixed scenes. Furthermore, the weighted random sampler assigns higher sampling weights to patches exhibiting larger depth variances, thereby compelling the model to focus on geometrically complex boundaries during optimization.

\subsection{(Three-Layer Angular Signal Decomposition)}

Conventional light field depth estimation methods predominantly rely on the Lambertian reflectance assumption [7], where the angular signal of a spatial point remains constant or forms linear structures in Epipolar Plane Images (EPIs) [7]. However, in real-world scenes, non-Lambertian materials such as specular highlights and translucent surfaces violate this assumption, causing severe distortions in the angular domain. Existing deep learning approaches typically process these distorted angular patches using black-box convolutional networks, which implicitly entangle illumination, geometry, and material properties, leading to depth ambiguities in complex regions. To address this physical boundary, we introduce the Three-Layer Angular Signal Decomposition module. Different from previous works that directly feed raw angular patches into feature extractors, our module explicitly decouples the discrete angular samples into illumination, disparity, and material components based on a physical lighting model, providing physically meaningful priors for the subsequent dual-branch network.

Let the input angular patch for a specific spatial pixel be denoted as $I(u,v) \in \mathbb{R}^{9 \times 9 \times C}$, where $(u,v)$ represents the discrete angular coordinates within the $9 \times 9$ micro-lens array, and $C$ is the number of color channels. According to the physical image formation model under varying viewpoints, the observed angular intensity can be approximated by a first-order Taylor expansion around the center viewpoint, combined with a residual term capturing high-frequency material reflections. We formulate the angular signal decomposition as:
\begin{equation}
\label{eq:eq6}
I(u,v) = DC + a \cdot u + b \cdot v + \varepsilon(u,v)
\end{equation}
where $DC$ represents the ambient illumination component, $a$ and $b$ are the linear coefficients capturing the disparity-induced EPI slopes along the horizontal and vertical angular dimensions, respectively, and $\varepsilon(u,v)$ is the residual term encoding material-specific high-frequency reflections. By explicitly isolating $\varepsilon(u,v)$, the network can effectively separate specular highlights from the underlying geometric structures, providing robust priors for the subsequent non-Lambertian branch.